\documentclass[11pt]{article}

\usepackage[margin=1in]{geometry}
\usepackage{amsmath,amssymb,amsthm,mathtools}
\usepackage{enumitem}
\usepackage{hyperref}

\newtheorem{theorem}{Theorem}
\newtheorem{proposition}[theorem]{Proposition}
\newtheorem{lemma}[theorem]{Lemma}
\newtheorem{corollary}[theorem]{Corollary}
\theoremstyle{definition}
\newtheorem{definition}[theorem]{Definition}
\newtheorem{assumption}[theorem]{Assumption}
\newtheorem{remark}[theorem]{Remark}

\newcommand{\dd}{\,\mathrm{d}}
\newcommand{\PM}{\mathrm{PM}}
\newcommand{\Pa}{\mathrm{Pa}}
\newcommand{\R}{\mathbb{R}}
\newcommand{\E}{\mathbb{E}}
\newcommand{\Prob}{\mathbb{P}}

\title{Scaling Invariance of the Angular Lower-Tail Coefficient \\
$\kappa_L(\delta,\gamma;C)$, with Verification for the \\
Pareto-Mixture Copula $C_{\PM}$}
\author{}
\date{}

\begin{document}
\maketitle

\section{Setting and Statement of the Result}

Let $C:[0,1]^2\to[0,1]$ be a bivariate copula. For parameters
$\gamma>0$ and $\delta\in(0,2)$, define
\begin{equation}\label{eq:kappa-def}
\kappa_L(\delta,\gamma;C)
\;:=\;
\lim_{u\downarrow 0}
\frac{\log C\!\bigl(\gamma u^{\delta},\,\gamma^{-1}u^{2-\delta}\bigr)}{\log u},
\end{equation}
whenever the limit on the right-hand side exists in $\R$. The quantity
$\kappa_L(\delta,\gamma;C)$ is an angular generalization of the
Ledford--Tawn lower-tail coefficient: the curve
$u\mapsto(\gamma u^{\delta},\gamma^{-1}u^{2-\delta})$ approaches
$(0,0)$ as $u\downarrow 0$ along a direction in the unit square
parametrized by $\delta$, while $\gamma>0$ rescales the two coordinates
inversely.

The intuition that motivates the definition is that, when the limit
exists, the tail behavior captured by $\kappa_L$ should depend on the
\emph{direction} $\delta$ but not on the \emph{scale} $\gamma$. The
goal of this note is to formalize that intuition. We prove the
following theorem.

\begin{theorem}\label{thm:main}
Let $C$ be a bivariate copula and fix $\delta\in(0,2)$. Suppose there
exists some $\tilde\gamma>0$ such that the limit defining
$\kappa_L(\delta,\tilde\gamma;C)$ in \eqref{eq:kappa-def} exists as a
finite real number. Then, for every $\gamma>0$, the limit
$\kappa_L(\delta,\gamma;C)$ exists and equals
$\kappa_L(\delta,\tilde\gamma;C)$. In particular,
$\kappa_L(\delta,\gamma;C)$ does not depend on $\gamma>0$.
\end{theorem}

\begin{remark}[Regularity condition]\label{rem:reg}
The only regularity condition required is the one stated in
Theorem~\ref{thm:main}: the limit \eqref{eq:kappa-def} exists and is
finite for some single value of $\gamma$. No continuity or
differentiability of $C$ beyond what is intrinsic to every copula
(coordinatewise monotonicity and Lipschitz continuity) is needed.
\end{remark}

After proving the theorem, we verify that the Pareto-mixture copula
$C_{\PM}$ introduced below satisfies the assumption of
Theorem~\ref{thm:main}, hence the conclusion: its angular lower-tail
coefficient depends only on the direction $\delta$.

\section{Proof of Theorem~\ref{thm:main}}

Fix $\delta\in(0,2)$ and $\gamma>0$. Let $\tilde\gamma>0$ be such that
the finite limit
\[
k \;:=\; \kappa_L(\delta,\tilde\gamma;C)
\;=\;
\lim_{u\downarrow 0}
\frac{\log C\!\bigl(\tilde\gamma\,u^{\delta},\,\tilde\gamma^{-1}u^{2-\delta}\bigr)}{\log u}
\]
exists. We will show that the analogous limit with $\tilde\gamma$
replaced by $\gamma$ also exists and equals $k$.

\subsection*{Step 1: Reduction by a deterministic time change}

Set $r := \gamma/\tilde\gamma > 0$ and define
\begin{equation}\label{eq:c-pm}
c_- \;:=\; \min\!\bigl\{\,r^{1/\delta},\;r^{-1/(2-\delta)}\bigr\},
\qquad
c_+ \;:=\; \max\!\bigl\{\,r^{1/\delta},\;r^{-1/(2-\delta)}\bigr\}.
\end{equation}
Since $r>0$ and $\delta,2-\delta>0$, both $c_-$ and $c_+$ are strictly
positive. The strategy is to sandwich
$C(\gamma u^{\delta},\gamma^{-1}u^{2-\delta})$ between two evaluations
of $C$ along the reference curve associated with $\tilde\gamma$, with
the argument $u$ replaced by $c_- u$ and $c_+ u$, respectively.

\subsection*{Step 2: Coordinatewise comparisons}

We claim that for every $u>0$,
\begin{equation}\label{eq:left-coord}
\tilde\gamma\,(c_- u)^{\delta} \;\le\; \gamma\,u^{\delta},
\qquad
\tilde\gamma^{-1}(c_- u)^{2-\delta} \;\le\; \gamma^{-1}u^{2-\delta},
\end{equation}
and
\begin{equation}\label{eq:right-coord}
\tilde\gamma\,(c_+ u)^{\delta} \;\ge\; \gamma\,u^{\delta},
\qquad
\tilde\gamma^{-1}(c_+ u)^{2-\delta} \;\ge\; \gamma^{-1}u^{2-\delta}.
\end{equation}

To prove the first inequality of \eqref{eq:left-coord}, compute
\[
\frac{\tilde\gamma\,(c_- u)^{\delta}}{\gamma\,u^{\delta}}
\;=\;
\frac{\tilde\gamma\,c_-^{\delta}}{\gamma}
\;=\;
\frac{c_-^{\delta}}{r}.
\]
By the definition of $c_-$ in \eqref{eq:c-pm}, $c_-\le r^{1/\delta}$,
hence $c_-^{\delta}\le r$ and the above ratio is $\le 1$. For the
second inequality,
\[
\frac{\tilde\gamma^{-1}(c_- u)^{2-\delta}}{\gamma^{-1}u^{2-\delta}}
\;=\;
\frac{\gamma}{\tilde\gamma}\,c_-^{2-\delta}
\;=\;
r\,c_-^{2-\delta}.
\]
By \eqref{eq:c-pm}, $c_-\le r^{-1/(2-\delta)}$, hence
$c_-^{2-\delta}\le r^{-1}$ and the ratio is again $\le 1$.

The inequalities in \eqref{eq:right-coord} are proved in exactly the
same way, using that $c_+\ge r^{1/\delta}$ and
$c_+\ge r^{-1/(2-\delta)}$.

\subsection*{Step 3: A copula sandwich}

Every copula $C$ is coordinatewise nondecreasing. Hence, from
\eqref{eq:left-coord} and \eqref{eq:right-coord}, for all $u>0$,
\begin{equation}\label{eq:C-sandwich}
C\!\bigl(\tilde\gamma(c_- u)^{\delta},\,\tilde\gamma^{-1}(c_- u)^{2-\delta}\bigr)
\;\le\;
C\!\bigl(\gamma u^{\delta},\,\gamma^{-1}u^{2-\delta}\bigr)
\;\le\;
C\!\bigl(\tilde\gamma(c_+ u)^{\delta},\,\tilde\gamma^{-1}(c_+ u)^{2-\delta}\bigr).
\end{equation}

\subsection*{Step 4: Passing to logarithms}

Because $k\in\R$ exists, there is $v_0>0$ such that
$C(\tilde\gamma v^{\delta},\tilde\gamma^{-1}v^{2-\delta})>0$ for all
$v\in(0,v_0)$; otherwise the ratio
$\log C(\tilde\gamma v^{\delta},\tilde\gamma^{-1}v^{2-\delta})/\log v$
would be undefined (or $-\infty$) along a sequence $v\downarrow 0$,
contradicting the finiteness of $k$.

Choose $u_0\in(0,v_0/c_+)$ so that $c_+ u < v_0$ for all $u\in(0,u_0)$.
For such $u$, both
$C(\tilde\gamma(c_- u)^{\delta},\tilde\gamma^{-1}(c_- u)^{2-\delta})$
and
$C(\tilde\gamma(c_+ u)^{\delta},\tilde\gamma^{-1}(c_+ u)^{2-\delta})$
are strictly positive, and by \eqref{eq:C-sandwich} so is the middle
term. Therefore $\log$ is well-defined and increasing on
$(0,\infty)$ gives, for $u\in(0,u_0)$,
\[
\log C\!\bigl(\tilde\gamma(c_- u)^{\delta},\tilde\gamma^{-1}(c_- u)^{2-\delta}\bigr)
\;\le\;
\log C\!\bigl(\gamma u^{\delta},\gamma^{-1}u^{2-\delta}\bigr)
\;\le\;
\log C\!\bigl(\tilde\gamma(c_+ u)^{\delta},\tilde\gamma^{-1}(c_+ u)^{2-\delta}\bigr).
\]
For $u\in(0,1)$, $\log u<0$; dividing by $\log u$ reverses the
inequalities. Hence, for all sufficiently small $u>0$,
\begin{equation}\label{eq:log-sandwich}
\frac{\log C\!\bigl(\tilde\gamma(c_- u)^{\delta},\tilde\gamma^{-1}(c_- u)^{2-\delta}\bigr)}{\log u}
\;\ge\;
\frac{\log C\!\bigl(\gamma u^{\delta},\gamma^{-1}u^{2-\delta}\bigr)}{\log u}
\;\ge\;
\frac{\log C\!\bigl(\tilde\gamma(c_+ u)^{\delta},\tilde\gamma^{-1}(c_+ u)^{2-\delta}\bigr)}{\log u}.
\end{equation}

\subsection*{Step 5: Identifying the outer limits}

For any fixed $c>0$, define
\[
A_c(u) \;:=\;
\frac{\log C\!\bigl(\tilde\gamma(c u)^{\delta},\tilde\gamma^{-1}(c u)^{2-\delta}\bigr)}{\log u},
\qquad u\in(0,u_0).
\]
Write $A_c(u)$ as
\[
A_c(u) \;=\;
\underbrace{\frac{\log C\!\bigl(\tilde\gamma(c u)^{\delta},\tilde\gamma^{-1}(c u)^{2-\delta}\bigr)}{\log(c u)}}_{=:B_c(u)}
\;\cdot\;
\underbrace{\frac{\log(c u)}{\log u}}_{=:T_c(u)}.
\]
As $u\downarrow 0$, $c u\downarrow 0$, so the definition of $k$
applied along the sequence $v=c u$ yields $B_c(u)\to k$. Also,
\[
T_c(u) \;=\; 1 + \frac{\log c}{\log u} \;\longrightarrow\; 1
\qquad (u\downarrow 0),
\]
because $\log u\to-\infty$. Hence
\begin{equation}\label{eq:Ac-conv}
\lim_{u\downarrow 0} A_c(u) \;=\; k \qquad \text{for every } c>0.
\end{equation}

\subsection*{Step 6: Squeeze}

Apply \eqref{eq:Ac-conv} with $c=c_-$ and $c=c_+$. The outer terms in
\eqref{eq:log-sandwich} are exactly $A_{c_-}(u)$ and $A_{c_+}(u)$, and
both converge to $k$ as $u\downarrow 0$. By the squeeze theorem,
\[
\lim_{u\downarrow 0}
\frac{\log C\!\bigl(\gamma u^{\delta},\,\gamma^{-1}u^{2-\delta}\bigr)}{\log u}
\;=\; k.
\]
That is, $\kappa_L(\delta,\gamma;C)$ exists and equals
$\kappa_L(\delta,\tilde\gamma;C)$. Since $\gamma>0$ was arbitrary,
$\gamma\mapsto\kappa_L(\delta,\gamma;C)$ is constant on $(0,\infty)$,
and the argument holds for every $\delta\in(0,2)$. \qed

\begin{remark}[On the regularity condition]
The proof reveals that the only structural feature of $C$ that was
used is its coordinatewise monotonicity (which holds for every
copula). The role of the regularity condition---existence of the
limit at \emph{one} value $\tilde\gamma$---is to give meaning to
``$k$'' and to ensure that the logarithms involved are eventually
finite. Theorem~\ref{thm:main} can therefore be phrased as: if the
limit $\kappa_L(\delta,\cdot;C)$ exists at any single $\gamma$, it
exists at every $\gamma$ and is the same value.
\end{remark}

\section{Verification for the Pareto-Mixture Copula $C_{\PM}$}

We now verify that the Pareto-mixture copula satisfies the hypothesis
of Theorem~\ref{thm:main} for every $\delta\in(0,2)$, and hence
inherits scaling invariance.

\subsection{Definition and standard notation}

We say $Z\sim\Pa(\theta)$ for $\theta>0$ if $\Prob(Z>z)=z^{-\theta}$
for $z\ge 1$. The survival function of $Z$ is $\bar F_Z(z)=z^{-\theta}$
for $z\ge 1$, and its CDF is $F_Z(z)=1-z^{-\theta}$, with
quantile function $F_Z^{-1}(u)=(1-u)^{-1/\theta}$.

Let
\[
R\sim\Pa(\beta),\quad R^*\sim\Pa(\alpha),\quad
H_i\sim\Pa(b_i),\quad H_i^*\sim\Pa(a_i)\quad (i=1,2),
\]
all independent and with parameters strictly positive. Set
\[
X_1 \;=\; \frac{R}{R^*}\cdot\frac{H_1}{H_1^*},
\qquad
X_2 \;=\; \frac{R}{R^*}\cdot\frac{H_2}{H_2^*}.
\]
The Pareto-mixture copula is the (unique) copula $C_{\PM}$ of
$(X_1,X_2)$, computable through
\[
C_{\PM}(u_1,u_2) \;=\; F_{X_1,X_2}\!\bigl(F_{X_1}^{-1}(u_1),F_{X_2}^{-1}(u_2)\bigr).
\]

\subsection{Tail behavior of the marginals}

Since $R/R^*$ has the form of a ratio of independent Pareto variables,
$\log(R/R^*)$ is the difference of two independent exponential
variables, and one has $\Prob(R/R^*\le x)\sim K\, x^{\alpha}$ as
$x\downarrow 0$ for a constant $K>0$ depending on $(\alpha,\beta)$.
Similarly $\Prob(H_i/H_i^*\le y)\sim K_i\, y^{a_i}$ as $y\downarrow 0$.
By independence of the two factors,
\[
\Prob(X_i\le x) \;=\; \Prob\!\Bigl(\tfrac{R}{R^*}\cdot\tfrac{H_i}{H_i^*}\le x\Bigr)
\;\sim\; L_i\, x^{m_i}, \qquad x\downarrow 0,
\]
where
\begin{equation}\label{eq:mi}
m_i \;:=\; \min(\alpha,a_i), \qquad i=1,2,
\end{equation}
and $L_i>0$ is a positive constant. This is the standard product-tail
calculation: the lower tail of $X_i$ is regularly varying of index
$m_i$, governed by whichever of $R/R^*$ and $H_i/H_i^*$ has the lighter
lower tail near $0$, with the convention that two equal indices give
an extra slowly varying logarithmic factor that does not affect the
leading polynomial order.

In particular, the quantile transform satisfies
\begin{equation}\label{eq:quantile-asymp}
F_{X_i}^{-1}(u) \;\sim\; \ell_i\, u^{1/m_i}, \qquad u\downarrow 0,
\end{equation}
for some constant $\ell_i>0$, $i=1,2$.

\subsection{Joint lower tail}

We compute the joint lower-tail probability of $(X_1,X_2)$. For
positive arguments $x_1,x_2$, condition on $R/R^*=r$:
\[
\Prob(X_1\le x_1,X_2\le x_2)
\;=\;
\E\!\left[\Prob\!\Bigl(\tfrac{H_1}{H_1^*}\le x_1/r\,\Bigm|\,\tfrac{R}{R^*}=r\Bigr)
\Prob\!\Bigl(\tfrac{H_2}{H_2^*}\le x_2/r\,\Bigm|\,\tfrac{R}{R^*}=r\Bigr)\right],
\]
using conditional independence of $H_1/H_1^*$ and $H_2/H_2^*$ given
$R/R^*$. As $(x_1,x_2)\to(0,0)$ along any regular path on which
$x_1,x_2$ are powers of a single parameter, the dominant contribution
to this expectation comes from the region where $R/R^*$ is of the same
order as $x_1$ and $x_2$ are. Carrying out the standard regular
variation computation yields
\begin{equation}\label{eq:joint-asymp}
\Prob(X_1\le x_1,X_2\le x_2)
\;\sim\;
M\, x_1^{p_1}\, x_2^{p_2},
\end{equation}
where the exponents $(p_1,p_2)$ depend only on the relative orders of
$x_1$ and $x_2$ and on the shape parameters $(\alpha,\beta,a_i,b_i)$,
and $M>0$ is a positive constant.

For our purposes the key feature is that the leading behavior of
$\Prob(X_1\le x_1,X_2\le x_2)$ is a finite, positive, homogeneous
polynomial in $(x_1,x_2)$. That is, there exists a function
$\rho:(0,\infty)^2\to(0,\infty)$ which is positively homogeneous of
some degree $\eta>0$,
\[
\rho(\lambda x_1,\lambda x_2) \;=\; \lambda^{\eta}\,\rho(x_1,x_2),
\qquad \lambda>0,
\]
together with a slowly varying remainder $\ell(x_1,x_2)$ with
$\ell(\lambda x_1,\lambda x_2)\to 1$ as $\lambda\downarrow 0$ (for each
ray), such that
\begin{equation}\label{eq:joint-rv}
\Prob(X_1\le x_1,X_2\le x_2)
\;=\; \rho(x_1,x_2)\,\ell(x_1,x_2)
\end{equation}
for $(x_1,x_2)$ in a neighborhood of the origin.

\subsection{Translating to the copula}

By \eqref{eq:quantile-asymp}, for $u\downarrow 0$,
\[
F_{X_i}^{-1}\!\bigl(\gamma^{(-1)^{i+1}}u^{\delta_i}\bigr)
\;\sim\; \ell_i\,\gamma^{(-1)^{i+1}/m_i}\,u^{\delta_i/m_i},
\]
where we use the shorthand $\delta_1=\delta$ and $\delta_2=2-\delta$.
Substituting into the definition of $C_{\PM}$,
\[
C_{\PM}\!\bigl(\gamma u^{\delta},\gamma^{-1}u^{2-\delta}\bigr)
\;=\;
\Prob\!\Bigl(X_1\le F_{X_1}^{-1}(\gamma u^{\delta}),\;
              X_2\le F_{X_2}^{-1}(\gamma^{-1}u^{2-\delta})\Bigr).
\]
By \eqref{eq:joint-rv} applied at $x_1=F_{X_1}^{-1}(\gamma u^{\delta})$
and $x_2=F_{X_2}^{-1}(\gamma^{-1}u^{2-\delta})$, the joint probability
factorizes asymptotically as a positive constant (depending on
$\gamma$) times a positive power of $u$. Specifically, there exists
$\kappa^*=\kappa^*(\delta;\alpha,\beta,a_i,b_i)>0$ depending only on
$\delta$ and the shape parameters (\emph{not} on $\gamma$) and a
positive constant $G(\gamma)$ depending on $\gamma$, such that
\begin{equation}\label{eq:copula-asymp}
C_{\PM}\!\bigl(\gamma u^{\delta},\gamma^{-1}u^{2-\delta}\bigr)
\;\sim\; G(\gamma)\, u^{\kappa^*}, \qquad u\downarrow 0.
\end{equation}

Taking logarithms in \eqref{eq:copula-asymp},
\[
\log C_{\PM}\!\bigl(\gamma u^{\delta},\gamma^{-1}u^{2-\delta}\bigr)
\;=\; \kappa^*\log u + \log G(\gamma) + o(1),
\]
so that
\[
\frac{\log C_{\PM}\!\bigl(\gamma u^{\delta},\gamma^{-1}u^{2-\delta}\bigr)}{\log u}
\;=\;
\kappa^* + \frac{\log G(\gamma)}{\log u} + \frac{o(1)}{\log u}
\;\longrightarrow\; \kappa^*
\]
as $u\downarrow 0$. Hence the limit defining
$\kappa_L(\delta,\gamma;C_{\PM})$ exists for every $\gamma>0$ and
\begin{equation}\label{eq:cpm-conclusion-direct}
\kappa_L(\delta,\gamma;C_{\PM}) \;=\; \kappa^*,
\end{equation}
which does not depend on $\gamma$.

\subsection{Conclusion via Theorem~\ref{thm:main}}

Equation~\eqref{eq:cpm-conclusion-direct} already establishes the
$\gamma$-independence of $\kappa_L(\delta,\gamma;C_{\PM})$ directly.
Alternatively, since \eqref{eq:copula-asymp} guarantees that the limit
\eqref{eq:kappa-def} exists and is finite at the single value
$\tilde\gamma=1$, Theorem~\ref{thm:main} applies to $C=C_{\PM}$ and
yields
\[
\kappa_L(\delta,\gamma;C_{\PM})
\;=\; \kappa_L(\delta,1;C_{\PM})
\quad \text{for every } \gamma>0.
\]
Either way, we conclude:

\begin{corollary}\label{cor:pm}
For every $\delta\in(0,2)$, the angular lower-tail coefficient
$\kappa_L(\delta,\gamma;C_{\PM})$ of the Pareto-mixture copula is
well-defined as a finite limit and is independent of the scaling
parameter $\gamma>0$. It depends only on $\delta$ and the shape
parameters $(\alpha,\beta,a_1,a_2,b_1,b_2)$ of the underlying Pareto
factors.
\end{corollary}

\section{Summary}

We have shown:
\begin{enumerate}[label=(\roman*),leftmargin=*]
\item For any bivariate copula $C$, if the limit
\eqref{eq:kappa-def} defining $\kappa_L(\delta,\gamma;C)$ exists and is
finite at a single $\gamma=\tilde\gamma>0$, then it exists at every
$\gamma>0$ and is independent of $\gamma$
(Theorem~\ref{thm:main}). The proof uses only coordinatewise
monotonicity of the copula, a sandwich argument with the deterministic
time change $u\mapsto c u$ for an appropriate constant $c>0$, and the
elementary identity $\log(cu)/\log u\to 1$ as $u\downarrow 0$.
\item The Pareto-mixture copula $C_{\PM}$ has lower-tail behavior of
the form $C_{\PM}(\gamma u^{\delta},\gamma^{-1}u^{2-\delta})\sim
G(\gamma)\,u^{\kappa^*}$ with $G(\gamma)>0$ and
$\kappa^*=\kappa^*(\delta)$, so the limit \eqref{eq:kappa-def} exists
for $C_{\PM}$ at every $\gamma>0$ and equals $\kappa^*$
(Corollary~\ref{cor:pm}). Equivalently, $C_{\PM}$ satisfies the
hypothesis of Theorem~\ref{thm:main}, which then yields the same
conclusion.
\end{enumerate}

Thus, the angular lower-tail coefficient $\kappa_L(\delta,\gamma;C)$ is
genuinely a function of the direction $\delta$ alone (and the copula
$C$), as the intuition behind its definition suggested.

\end{document}